\par\SourceLeaf{145}\noindent
Um diese vollkommen verständlich zu machen, bin ich genöthigt, einige Entwicklungen vorauszuschicken, in welchen mir zu folgen mancher vielleicht beschwerlich finden könnte, wenn von vorneherein noch nicht abzusehen ist, auf was sie hinauslaufen werden. Es wird deshalb, deucht mir, angemessen sein, wenn ich die vornehmsten Ausstellungen, welche das \textsc{Kritter}’sche Gutachten und die daraus gezogenen Folgerungen treffen, gleich hier an die Spitze stelle.
\begin{list}{}{\setlength{\leftmargin}{2em}\setlength{\labelwidth}{1.5em}\setlength{\labelsep}{.5em}\setlength{\itemsep}{0pt}\setlength{\parsep}{0pt}}
\item[1)] Die Anwendbarkeit des Verhältnisses von 1 Witwe gegen 2 stehende Ehen, auf die hiesige Professoren-Witwenkasse, ist nicht sicher; es ist vielmehr, wie schon oben S. [133] Z. [5] bemerkt ist, wahrscheinlich, dass nach den Verhältnissen dieser Kasse etwas mehr an Witwen gerechnet werden müsse.
\item[2)] \textsc{Kritter}’s Behauptung [Z. 10 des Abdrucks] ist, auch wenn sie unabhängig von diesem vorgesetzten Verhältnisse \(2:1\) vorgetragen wird, nemlich man könne annehmen, dass ein Durchschnitt von 10 Jahren immer schon hinreiche, das wahre für die Umstände der Kasse gültige Verhältniss sehr nahe anzugeben, ist durchaus falsch, und die Unrichtigkeit dieses Satzes ist auch ohne Berufung auf Erfahrungen, schon aus theoretischen Gründen nachzuweisen.
\item[3)] Weit wichtiger als diese beiden Ausstellungen ist der Umstand, dass P. das \textsc{Kritter}’sche Gutachten falsch ausgelegt hat, indem das, was P. unter Maximum der Witwenzahl verstand, und das, was man mit diesem Worte bezeichnet, wenn ein bestimmtes Normalverhältniss zwischen stehenden Ehen und Witwen aufgestellt wird, und was auch in \textsc{Kritter}’s Gutachten eigentlich gemeint ist,
\begin{center}\textit{zwei sehr verschiedene Dinge sind.}\end{center}
Hierzu kommt noch der eben so wichtige Umstand
\item[4)] dass die numerischen Resultate, die für eine Gesellschaft von einem gewissen sich immer nahe gleich bleibenden Umfang zulässig waren, wesentlich abgeändert werden müssen, wenn dieser Umfang sich bedeutend erweitert, wie diess schon oben ausführlich abgehandelt ist.
\end{list}
\par\medskip\centerline{\rule{.25\linewidth}{.3pt}}\medskip
\noindent\hspace*{\parindent}Ich fange an mit der (fingirten) Annahme, dass durch das Zusammentreten einer sehr grossen Anzahl von Ehepaaren aus den verschiedensten Alterstufen eine Gesellschaft gebildet worden sei. Alljährlich wird eine Anzahl von Ehen durch den Tod des einen oder des andern Theils getrennt werden: dieser Abgang werde dadurch ersetzt, dass jährlich eine bestimmte Zahl neuer Ehepaare hinzutritt, so viele, dass im Ganzen der Bestand der Gesellschaft ungeändert bleibe. Von den im Laufe eines Jahres durch den Tod des Ehemannes entstehenden Witwen wird, da die Gesellschaft als sehr gross vorausgesetzt wird, ein verhältnissmässiger sehr kleiner Theil schon während desselben Jahres wieder absterben, wodurch mithin die Anzahl der am Ende des Jahres wirklich vorhandenen Witwen etwas modificirt wird. Ungefähr eben so viele neue Witwen werden am Ende des zweiten Jahres hinzugekommen, dagegen aber von den aus dem ersten Jahre herrührenden Witwen ein Theil schon wieder verstorben sein, so dass am Ende des zweiten Jahres die Anzahl der Witwen nicht ganz doppelt so gross sein wird, als am Ende des ersten. Am Ende des dritten Jahres sind wieder ungefähr eben so viele neue Witwen hinzugekommen, dagegen wird der Bestand, welcher zu Anfang des dritten Jahres vorhanden war, einen fast doppelt so grossen Abgang erlitten haben, als die Bilanz des zweiten Jahres ergeben hatte. Man sieht, dass auf diese Weise die Zahl der Witwen zwar fortwährend wächst, aber immer langsamer, bis sie zuletzt so gross geworden ist, dass der einjährige Abgang durch den Tod der Witwen, dem Zugang durch Absterben von Ehemännern aus der Gesellschaft, das Gleichgewicht hält: dann wird also der Beharrungszustand eintreten, indem die Witwenzahl ihr Maximum erreicht hat, und von da an ungeändert bleibt.

Fügen wir den obigen Voraussetzungen noch die bei, dass sowohl die Anzahl der in jedem Jahre
